% Options for packages loaded elsewhere
\PassOptionsToPackage{unicode}{hyperref}
\PassOptionsToPackage{hyphens}{url}
\PassOptionsToPackage{dvipsnames,svgnames,x11names}{xcolor}
%
\documentclass[
  11pt,
]{article}
\usepackage{amsmath,amssymb}
\usepackage{iftex}
\ifPDFTeX
  \usepackage[T1]{fontenc}
  \usepackage[utf8]{inputenc}
  \usepackage{textcomp} % provide euro and other symbols
\else % if luatex or xetex
  \usepackage{unicode-math} % this also loads fontspec
  \defaultfontfeatures{Scale=MatchLowercase}
  \defaultfontfeatures[\rmfamily]{Ligatures=TeX,Scale=1}
\fi
\usepackage{lmodern}
\ifPDFTeX\else
  % xetex/luatex font selection
\fi
% Use upquote if available, for straight quotes in verbatim environments
\IfFileExists{upquote.sty}{\usepackage{upquote}}{}
\IfFileExists{microtype.sty}{% use microtype if available
  \usepackage[]{microtype}
  \UseMicrotypeSet[protrusion]{basicmath} % disable protrusion for tt fonts
}{}
\makeatletter
\@ifundefined{KOMAClassName}{% if non-KOMA class
  \IfFileExists{parskip.sty}{%
    \usepackage{parskip}
  }{% else
    \setlength{\parindent}{0pt}
    \setlength{\parskip}{6pt plus 2pt minus 1pt}}
}{% if KOMA class
  \KOMAoptions{parskip=half}}
\makeatother
\usepackage{xcolor}
\usepackage[margin=1in]{geometry}
\usepackage{longtable,booktabs,array}
\usepackage{calc} % for calculating minipage widths
% Correct order of tables after \paragraph or \subparagraph
\usepackage{etoolbox}
\makeatletter
\patchcmd\longtable{\par}{\if@noskipsec\mbox{}\fi\par}{}{}
\makeatother
% Allow footnotes in longtable head/foot
\IfFileExists{footnotehyper.sty}{\usepackage{footnotehyper}}{\usepackage{footnote}}
\makesavenoteenv{longtable}
\setlength{\emergencystretch}{3em} % prevent overfull lines
\providecommand{\tightlist}{%
  \setlength{\itemsep}{0pt}\setlength{\parskip}{0pt}}
\setcounter{secnumdepth}{-\maxdimen} % remove section numbering
\ifLuaTeX
  \usepackage{selnolig}  % disable illegal ligatures
\fi
\usepackage{bookmark}
\IfFileExists{xurl.sty}{\usepackage{xurl}}{} % add URL line breaks if available
\urlstyle{same}
\hypersetup{
  pdftitle={Entropy slopes in Kähler classes: calibrated examples and an audit of recovery},
  pdfauthor={Research draft prepared in the supplied workspace},
  colorlinks=true,
  linkcolor={Maroon},
  filecolor={Maroon},
  citecolor={Blue},
  urlcolor={Blue},
  pdfcreator={LaTeX via pandoc}}

\title{Entropy slopes in Kähler classes: calibrated examples and an
audit of recovery}
\author{Research draft prepared in the supplied workspace}
\date{6 October 2026}

\begin{document}
\maketitle

\textbf{Status: PARTIAL.} This is a standalone mathematical paper about
the unchanged P19.1 question. It does not certify entropy recovery for
every maximal finite-energy ray in an arbitrary transcendental Kähler
class. In particular it does not inherit a proof from Own178, the
earlier scoped counterexample verdict, or an abstract of a primary
source. The results below are explicit analytic examples, a
normalization calculation, a local counterexample to a step in an
imported proof, a log-canonical divisor inequality, and precisely
conditional recovery and ray-transfer statements. None is counted as a
complete solution of P19.1. P19.2, P19.3 and P19.4 are unclaimed.

\section{1. Conventions, hypotheses and background
inputs}\label{conventions-hypotheses-and-background-inputs}

Let \(X\) be a connected compact smooth Kähler manifold of dimension
\(n\geq1\), let \(\omega\) be a Kähler form, and put
\(\alpha=[\omega]\), \(V=\int_X\omega^n\), and \(\mu_0=\omega^n/V\).
Throughout the main paper, \[
 dd^c_C=\frac{i}{2\pi}\partial\bar\partial,\qquad
 t=-\log|\zeta|^2,\qquad
 dd^c_C\log|f|^2=[\operatorname{div}(f)].                 \tag{1.1}
\] We use bounded Bedford--Taylor products and their finite-energy
extensions. Set \[
 E(u)=\frac1{(n+1)V}\sum_{j=0}^n\int_Xu\,\omega_u^j\wedge\omega^{n-j},
 \qquad \operatorname{MA}(u)=V^{-1}\omega_u^n,
 \qquad \omega_u=\omega+dd^c_Cu.                         \tag{1.2}
\] The reference energy is \(E(0)=0\). Probability entropy is \[
 H(u)=\operatorname{Ent}(\operatorname{MA}(u)\mid\mu_0)
 =\int\rho\log\rho\,d\mu_0                              \tag{1.3}
\] when \(\operatorname{MA}(u)=\rho\mu_0\), and is \(+\infty\)
otherwise. It is nonnegative. Changing the smooth positive reference
probability adds a uniformly bounded term, so it does not change an
entropy slope.

An actual weak geodesic ray is a family of finite-energy potentials,
strongly continuous also at time zero, whose restriction to every finite
time interval is the canonical pluripotential weak geodesic. Bounded
means bounded at each finite time. It does not mean uniformly bounded
for all time. The explicit rays constructed below are bounded at every
finite time and have a smooth zero starting point, except the stationary
example in Section 4.

For a smooth dominating \(\mathbf C^*\)-equivariant configuration
\(\rho:\mathcal X\to X\times\mathbf P^1\), write
\(B=\mathcal X_0=\sum b_iE_i\) and \(A=\rho_X^*\alpha\). All NA model
divisors and model modifications in this paper are equivariant. A
vertical invariant real Cartier divisor \(D\) gives the model function
\[
 f_D(v_{E_i})=\frac{\operatorname{ord}_{E_i}D}{b_i}.
                                                               \tag{1.4}
\] The Gauss valuation has time value one. A compatible total potential
satisfies \(U\circ\rho+\log|f_D|^2=O(1)\), interpreted locally and
linearly for real divisors. Thus its divisor coefficient in (1.1) is
\(-D\). More generally the ordinary datum on a divisorial valuation is
minus the generic coefficient of \(U\) on its vertical Gauss divisor
divided by \(b_i\), after an affine time shift if needed to make \(U\)
locally bounded above. This prescription retains supports of
semivaluations and is not a new roof metric or an assigned tropical
datum. Formula (1.5) below concerns Gauss divisors on equivariant
models; it is not a discrepancy identity for the restriction of every
nonequivariant total-space divisorial valuation.

Log discrepancy is over \(X\). For a divisor \(F\) on a smooth
modification of \(X\),
\(A_X(c\operatorname{ord}_F)=c(1+\operatorname{ord}_FK_{Z/X})\), and
\(A_X(v_{\rm triv})=0\). For a vertical component on \(\mathcal X\), put
\(a_i=\operatorname{ord}_{E_i}K_{\mathcal X/(X\times\mathbf P^1)}\). Its
Gauss restriction has \[
 b_iA_X(v_{E_i})=a_i+1-b_i.
                                                               \tag{1.5}
\] The normalized model measure, for a relatively nef class
\(\Gamma=A+D\), is \[
 \operatorname{MA}^{\rm NA}(f_D)
 =V^{-1}\sum_i b_i(\Gamma^n\cdot E_i)\delta_{v_{E_i}},
 \qquad H^{\rm NA}(\phi)=\int A_X\,d\operatorname{MA}^{\rm NA}(\phi).
                                                               \tag{1.6}
\] For the explicit surface and product models below these are defining
model calculations, not uses of an entropy-slope theorem. For an
arbitrary PL obstacle, the positive-product envelope formulas are stated
separately as an external input in Section 7.

We use the following ordinary analytic background precisely in the
following forms. These are external foundations, not numerical findings.

\begin{itemize}
\tightlist
\item
  Locally bounded psh functions have Bedford--Taylor products; on a
  complex manifold their homogeneous Monge--Ampère equation gives the
  weak Dirichlet geodesic with prescribed bounded boundary values by the
  comparison principle. This is the Bedford--Taylor Dirichlet/comparison
  theory {[}BT{]}, applied only on a finite annulus with bounded
  potentials. For a \(C^{1,1}\) local potential, convolution produces
  smooth psh potentials converging locally uniformly, with uniformly
  bounded Hessians. At Lebesgue points those Hessians converge almost
  everywhere. Dominated convergence identifies their determinant
  measures, while local uniform continuity of Bedford--Taylor products
  identifies the same limit with the product of the original potential.
  This proves the determinant identification used below.
\item
  Bounded potentials lie in \(E^1\). Aubin--Mabuchi energy is affine on
  actual finite-energy geodesics. Its supporting inequalities are \[
  \int(u-v)\operatorname{MA}(u)\leq E(u)-E(v)
     \leq\int(u-v)\operatorname{MA}(v).                 \tag{1.7}
  \] The normalized \(d_1\) metric satisfies \(d_1(u,v)=E(v)-E(u)\) for
  \(u\leq v\), and in general \(d_1(u,v)=E(u)+E(v)-2E(P(u,v))\). These
  are the finite-energy metric and energy facts in Darvas {[}D{]} and
  Berman--Darvas--Lu {[}BDL{]}, Sections 2 and 4.2. Their use below is
  on \(E^1\) on a fixed compact Kähler manifold, with no rationality
  assumption.
\item
  The extended cohomological Mabuchi functional in (2.4) is \(d_1\)
  lower semicontinuous and convex on finite-energy geodesics. Smooth
  potentials approximate any \(E^1\) potential with simultaneous \(d_1\)
  and entropy convergence. The exact imports are {[}BDL,
  arXiv:1510.01260v2, Theorems 3.2 and 4.7{]}, with \(p=1\), twisting
  \(f=0\), and the change of potential scale detailed in Section 2.
  Their proofs use density truncation and complex Monge--Ampère
  solutions, then finite-time geodesic approximation. Entropy continuity
  on arbitrary strongly convergent sequences is not asserted.
\item
  A psh function on the complement of a closed pluripolar set, locally
  bounded above near that set, has a unique psh extension {[}Dem,
  Chapter I, Theorem 5.24{]}. A divisor is locally pluripolar by
  \(\log|f|\). For a closed positive \((1,1)\) current \(T\) and a prime
  divisor \(F\) on a smooth manifold, subtracting its generic
  coefficient gives \(T-\lambda[F]\geq0\) {[}Dem, Chapter III, Section
  8, Proposition 8.16 preceding the decomposition formula{]}. Locally
  the remainder has a psh potential, so divisorial subtraction gives
  \(q\leq\lambda\log|w|^2+O(1)\), with a smooth local background
  potential absorbed in the bound. Section 5 applies exactly these
  statements to the total potential after a smooth local background is
  added.
\end{itemize}

\section{2. Calibration of time, energy, Ricci curvature and
entropy}\label{calibration-of-time-energy-ricci-curvature-and-entropy}

\textbf{Proposition 2.1 (the vertical constant test).} On the product
model with \(D=cX_0\), the compatible ray in (1.1) is \(u_t=ct\). It has
datum \(c\), energy slope \(c\), and model energy \(c\). Keeping
\(dd^c_C\) but changing time to \(s=-\log|\zeta|\) gives the ray
\(u_s=2cs\) and energy slope \(2c\).

\textbf{Proof.} Compatibility means \(U=-c\log|\zeta|^2+O(1)\). For the
canonical constant ray the bounded remainder is zero. Adding a constant
to a potential adds that constant to (1.2), because each of the \(n+1\)
mixed products has mass \(V\). On \(X\times\mathbf P^1\), \(X_0^2=0\),
\(A^{n+1}=0\), and \(A^n\cdot X_0=V\). Therefore \[
 \frac{(A+cX_0)^{n+1}}{(n+1)V}=c.                         \tag{2.1}
\] The unique central component has weight one and discrepancy zero.
Formula (1.4) gives datum \(c\), and the two time claims follow. This
proves a genuine factor-of-two failure of the printed combination of
conventions in Own178. Zero entropy alone could not detect it.
\(\square\)

The Green--Riesz factor can be checked without test configurations. For
an \(S^1\)-invariant scalar function \(F\) on
\(\mathbf P^1\setminus\{0\}\) smooth away from zero, with no singularity
at infinity, Stokes gives, whenever the expressions are defined, \[
 \int_{\mathbf P^1\setminus\Delta_r}dd^c_CF
 =\tfrac12\frac{dF}{ds},\qquad s=-\log r.
                                                               \tag{2.2}
\] Indeed \(dd^c_CF=(4\pi)^{-1}\Delta F\,dA\), and the boundary flux is
\(-rF_r/2=F_s/2\). In time \(t=2s\) the right-hand side is \(dF/dt\).
Thus the energy push-forward argument in {[}SD, Theorem 4.9{]} requires
a factor two with that source's printed \(dd^c_C\) and
\(s=-\log|\zeta|\). Simply copying its displayed derivative equation
does not pass (2.1).

Define \[
 \operatorname{Ric}_C(\omega)=-dd^c_C\log\omega^n,\qquad
 \overline S_C=\frac nV\int_X\operatorname{Ric}_C(\omega)\wedge\omega^{n-1},
\] so \([\operatorname{Ric}_C]=c_1(X)\). With \[
 T_\chi(u)=V^{-1}\sum_{j=0}^{n-1}\int_Xu\,\chi\wedge\omega_u^j\wedge\omega^{n-1-j},
                                                               \tag{2.3}
\] the cohomological Mabuchi functional is \[
 M_C(u)=H(u)+T_{-\operatorname{Ric}_C(\omega)}(u)+\overline S_C E(u),
 \qquad M_C(0)=0.                                             \tag{2.4}
\] For a smooth variation \(\dot u\), differentiating the entropy and
integrating by parts gives \[
 \dot H=\frac nV\int_X\dot u\,
     (\operatorname{Ric}_C(\omega)-\operatorname{Ric}_C(\omega_u))
        \wedge\omega_u^{n-1}.
\] The derivative of \(T_\chi\) is
\(nV^{-1}\int\dot u\,\chi\wedge\omega_u^{n-1}\). Consequently
\(\dot M_C=\int\dot u(\overline S_C-S_C(\omega_u))\operatorname{MA}(u)\),
verifying the entropy coefficient in (2.4).

For comparison set \(dd^c_L=i\partial\bar\partial/\pi\) and keep the
same background form \(\omega\). The precise dictionary preserving the
same currents, class, total potential divisor weights and model datum is
\[
 v_s=\tfrac12u_{2s},\qquad s=-\log|\zeta|,
 \quad E_L(v_s)=\tfrac12E_C(u_{2s}),
 \quad H_L(v_s)=H_C(u_{2s}).                                \tag{2.5}
\] The cohomological Ricci form in this convention is
\(\operatorname{Ric}_{\rm coh}=-(1/2)dd^c_L\log\omega^n=\operatorname{Ric}_C\).
The correct cohomological Mabuchi expression is \[
 M_L=\tfrac12 H_L+T_{-\operatorname{Ric}_{\rm coh},L}
                      +\overline S_C E_L.
                                                               \tag{2.6}
\] Indeed \(M_L(v)=M_C(2v)/2\), so the two cohomological Mabuchi slopes
agree under (2.5), whereas the ordinary probability entropy slope in
\(s\) is twice its slope in \(t\). BHJ {[}BHJ, Section 1 and Lemma
3.11{]} use logarithmic metric weights, half probability entropy, and
volume \(e^{2\Psi}\); this agrees with (2.5)--(2.6). {[}MP, Definition
5.1.2{]} uses \(\log|f|\) and time \(-\log|\zeta|\). Its mixed and
entropy formulas must be transported, not copied as squared-log
formulas. We do not certify every printed scalar formula in either
source.

Finally {[}BDL{]} uses \(i\partial\bar\partial\). Its potential
\(q=u/(2\pi)\) has the same current as our \(u\). Its Ricci form is
\(2\pi\operatorname{Ric}_C\), its energy is \(E_C/(2\pi)\), and its
contracted-energy products give exactly (2.4). Its metric \(d_1\) is our
\(d_1/(2\pi)\). Thus its lower semicontinuity, convexity and entropy
approximation apply to (2.4) without a change of the probability
entropy.

\section{3. An explicit nonzero divisor entropy
calibration}\label{an-explicit-nonzero-divisor-entropy-calibration}

Take \(X=\mathbf P^1\) and \(\omega=dd^c_C\log(1+|z|^2)\), of mass one.
Put \[
 x=\log|z|^2,\quad h(x)=\log(1+e^x),\quad
 y=h'(x)=\frac{e^x}{1+e^x},\quad
 g(y)=y\log y+(1-y)\log(1-y).
\] Then \(\omega=dy\,d\theta/(2\pi)\). Fix \(0<\varepsilon<1\), write
\(x_\varepsilon=\log(\varepsilon/(1-\varepsilon))\), and define \[
 h_t(x)=\sup_{0\leq y\leq1}\{xy-g(y)-t(\varepsilon-y)_+\},
 \qquad u_t=h_t-h.                                        \tag{3.1}
\]

\textbf{Theorem 3.1.} The functions (3.1) form an actual bounded weak
geodesic ray from zero. In convention (1.1), \[
 -\varepsilon t\leq u_t\leq0,\qquad E(u_t)=-\varepsilon^2t/2,\qquad
 \lim_{t\to\infty}\frac{H(u_t)}t=\varepsilon.                        \tag{3.2}
\] Its ordinary datum is the relatively Kähler
deformation-to-the-normal-cone model \(D=-\varepsilon E\), where \(E\)
is the exceptional divisor of the blow-up of \((z,\zeta)=(0,0)\). Its
normalized NA measure is
\((1-\varepsilon)\delta_{v_{\rm triv}}+\varepsilon\delta_{\operatorname{ord}_0}\),
and its NA entropy is \(\varepsilon\). Thus the equality is calibrated
on a genuine positive-discrepancy divisor, not just on entropy zero.
This is a model special case, not a resolution of the residual
maximal-ray atom.

\textbf{Proof.} The explicit formula is \[
 h_t(x)=
 \begin{cases}
 h(x+t)-\varepsilon t,&x<x_\varepsilon-t,\\
 \varepsilon x-g(\varepsilon),&x_\varepsilon-t\leq x\leq x_\varepsilon,\\
 h(x),&x>x_\varepsilon.
 \end{cases}                                                \tag{3.3}
\] The values and first derivatives in \((x,t)\) agree on the
interfaces. It is convex jointly in \((x,t)\) because (3.1) is a
supremum of affine functions. Its real Hessian has rank at most one on
each region. On each compact finite-time cylinder it is \(C^{1,1}\),
including the interfaces. Near \(z=0\),
\(u_t=\log(1+e^t|z|^2)-\varepsilon t-\log(1+|z|^2)\), which extends
smoothly for finite time. Near infinity \(u_t=0\). The total toric
potential is therefore locally bounded psh on
\(\mathbf P^1\times\Delta^*\) and solves the homogeneous complex
Monge--Ampère equation: on \(\mathbf C^*\times\Delta^*\) the determinant
is zero almost everywhere; the \(C^{1,1}\) Bedford--Taylor
identification gives the equation; near \(z=0\) it follows from the
displayed smooth expression. On a finite annulus comparison and
uniqueness identify it with the weak geodesic with its boundary values.
Uniform continuity on finite-time compact sets gives strong continuity
at zero. The bounds follow immediately from
\(0\leq(\varepsilon-y)_+\leq \varepsilon\) in (3.1).

For energy, \(\partial_t h_t=-(\varepsilon-y_t)_+\) almost everywhere,
where \(y_t=\partial_xh_t\). The energy derivative is the MA integral of
\(\dot u_t\). On the left region \(y_t=h'(x+t)\) runs from zero to
\(\varepsilon\); the middle region has MA measure zero; the right region
contributes zero to this derivative. Hence \[
 \frac{d}{dt}E(u_t)=-\int_0^\varepsilon(\varepsilon-y)\,dy=-\varepsilon^2/2.
\] The energy derivative is legitimate for this locally Lipschitz path
of bounded potentials by bounded product approximation (or by directly
differentiating the curve energy
\(E(u)=\frac12\int u(\omega+\omega_u)\)). Since \(E(u_0)=0\), the energy
claim follows.

Let \(k(x)=h''(x)=e^x/(1+e^x)^2\). The MA density relative to \(\omega\)
is \(k(x+t)/k(x)\) on \(x<x_\varepsilon-t\), zero on the middle annulus,
and one on \(x>x_\varepsilon\). There is no interface atom, since
\(\partial_xh_t\) is continuous. Therefore, with \(q=x+t\), \[
 \begin{split}
 H(u_t)&=\int_{-\infty}^{x_\varepsilon}
      \log\frac{k(q)}{k(q-t)}\,k(q)\,dq\\
 &=\varepsilon t+2\int_{-\infty}^{x_\varepsilon}
       \{\log(1+e^{q-t})-\log(1+e^q)\}\,k(q)\,dq.
 \end{split}                                                \tag{3.4}
\] The second term is between \(-2\varepsilon\log(1+e^{x_\varepsilon})\)
and zero. In particular every finite-time entropy is finite and the
slope is \(\varepsilon\).

It remains to identify the actual total-space datum. Formula (3.3) shows
in the three regions, near \((0,0)\), that the total potential differs
by a bounded function from \(\varepsilon\log(|z|^2+|\zeta|^2)\). On the
region \(|z|^2\leq e^{x_\varepsilon}|\zeta|^2\) use the first formula;
on the middle region use \(h_t=\varepsilon x-g(\varepsilon)\) and
\(|z|^2\geq e^{x_\varepsilon}|\zeta|^2\); the final region is bounded
away from \(z=0\). Thus the ideal \((z,\zeta)\) gives coefficient
\(\varepsilon E\) after the blow-up, and
\(U+\log|f_{-\varepsilon E}|^2\) is bounded there. This identifies the
datum on every higher divisor too, by pullback of the bounded remainder.

On the blown-up surface let \(B\) denote the pullback of a fiber and
\(A\) the pullback of \([\omega]\). Then \[
 A^2=B^2=A\cdot E=B\cdot E=0,\quad A\cdot B=1,\quad E^2=-1.
\] The central fiber is \(B'+E\), with \(B'=B-E\) and both
multiplicities one. For \(\Gamma=A-\varepsilon E\), \[
 \Gamma\cdot B'=1-\varepsilon,\qquad\Gamma\cdot E=\varepsilon,
 \qquad\Gamma^2/2=-\varepsilon^2/2.                                 \tag{3.5}
\] The exceptional discrepancy over the smooth surface is \(a_E=1\), so
(1.5) gives \(A_X(v_E)=1\); the strict component has discrepancy zero
and trivial restriction. This proves the NA measure and entropy
assertions from (1.6). The configuration is relatively Kähler:
\(A+cB-\varepsilon E\) is Kähler for \(c>\varepsilon\) and
\(\varepsilon<1\). Here is a direct verification of that last fact. For
rational \(c,\varepsilon\), clear denominators to obtain integers
\(a,b,d>0\) with \(a,b>d\). The monomials of \(\mathcal O(a,b)\) whose
exponents satisfy \(i+j\geq d\) give sections of
\(\mathcal O(aA+bB-dE)\) on the blow-up. They have no common zero. On
the chart \(y=xr\) over the blown-up corner, division by the vertex
section \(x^d\) gives the two ratios \(x\) and \(r\) from \(x^{d+1}\)
and \(x^{d-1}y\). They recover the local coordinates. The other
exceptional chart gives \(y\) and \(x/y\) in the same way. At each of
the other three corners the vertex section and its two adjacent
monomials recover the ordinary affine coordinates. These five charts are
the nonvanishing loci of the five vertex sections and cover the blow-up.
On each locus the section map is an embedding because the two coordinate
ratios recover its point and tangent directions; points on different
loci are separated by the nonvanishing pattern and those coordinates.
Properness then gives a projective embedding. Pulling back a
Fubini--Study form and dividing by \(a\) proves the rational class is
Kähler. Any real \((c,\varepsilon)\) with the same strict inequalities
is a convex combination of finitely many rational pairs in that open
region. The corresponding positive combination of the constructed forms
proves the real case. Adding \(cB\) translates the ray by \(ct\), and
changes neither entropy nor the model masses. \(\square\)

\textbf{Product extension.} For any compact Kähler \((Y,\eta)\) of
dimension \(m\geq0\), the ray \(p_1^*u_t\) on \(\mathbf P^1\times Y\)
with form \(\Theta=p_1^*\omega+p_2^*\eta\) has exactly the energy and
entropy in (3.2). It is an actual product weak ray. To see the last
assertion, restrict a Perron competitor to each
\(\mathbf P^1\times\{y\}\) and its time annulus; the restriction is a
curve competitor or is identically \(-\infty\). The curve comparison
bounds it above by \(u_t\); the pullback supplies the reverse
inequality. The products are \[
 \Theta_{p_1^*u}^{m+1}=(m+1)\omega_u\wedge\eta^m,
 \qquad V_\Theta=(m+1)\int_Y\eta^m,
\] and hence the MA probabilities are a tensor product with the same
\(Y\) probability for the reference and the ray. The energy equality
follows from its derivative along \(s\mapsto p_1^*(su)\) and the value
at \(s=0\). The total model is the product of the surface blow-up with
\(Y\), so the same normalized weights and discrepancy values follow by
the identical intersection cancellation. This applies in transcendental
product classes and on nonprojective products whenever \(Y\) has those
properties. It is still a specific model ray.

\section{4. A bounded infinite-entropy starting
potential}\label{a-bounded-infinite-entropy-starting-potential}

The starting hypothesis cannot be silently removed even for bounded
starting potentials.

\textbf{Theorem 4.1.} There is a bounded \(\omega\)-psh function \(w\)
on \(\mathbf P^1\), with \(\omega\) as in Section 3, for which
\(H(w)=+\infty\). Its stationary actual weak ray \(w_t=w\) has ordinary
NA datum zero, \(H^{\rm NA}(0)=0\), energy slope zero, and
\(H(w_t)/t=+\infty\) for every \(t>0\). The construction also works on
\(\mathbf P^1\times Y\) in every product class of Section 3.

\textbf{Proof.} Use the coordinate \(y=|z|^2/(1+|z|^2)\) and the uniform
angular probability. Write \(Z=\sum_{j\geq1}j^{-2}\),
\(m_j=(2Zj^2)^{-1}\), and \(\ell_j=2^{-j-4}e^{-j^3}\). Choose disjoint
intervals \[
 I_j=\left[\tfrac13+2\sum_{i<j}\ell_i,
               \tfrac13+2\sum_{i<j}\ell_i+\ell_j\right].
\] All are inside \([1/3,1/2)\), since \(2\sum\ell_j<1/8<1/6\). Define
\[
 \rho(y)=\tfrac12+\sum_{j\geq1}\frac{m_j}{\ell_j}\mathbf1_{I_j}(y),
 \qquad \sigma=\rho\,\omega.
                                                               \tag{4.1}
\] Its mass is \(1/2+\sum m_j=1\). On each interval \(m_j/\ell_j>1\);
since \(r\log r\) is increasing for \(r\geq1\), its contribution to
entropy is at least \(m_j\log(m_j/\ell_j)\). The sum diverges because
its summands grow like a positive constant times \(j\). The negative
contribution on the complement is finite. Thus
\(\operatorname{Ent}(\sigma\mid\omega)=+\infty\).

Let \[
 G(z,a)=\log\frac{|z-a|^2}{(1+|z|^2)(1+|a|^2)},
 \qquad w(z)=\int G(z,a)\,d\sigma(a).                       \tag{4.2}
\] Then \(dd^c_CG(\cdot,a)=\delta_a-\omega\). The background half of
\(\sigma\) has constant finite Green potential, by the invariance of
\(\omega\) and the chordal kernel. Each spike is angularly uniform at
radii in a fixed compact interval \([r_-,r_+]\subset(0,\infty)\). The
elementary circle mean formula is \[
 \frac1{2\pi}\int_0^{2\pi}G(z,re^{i\theta})\,d\theta
 =\log\max\{|z|^2,r^2\}-\log(1+|z|^2)-\log(1+r^2).
\] The right-hand side is uniformly bounded for all \(z\in\mathbf P^1\)
and \(r\in[r_-,r_+]\). For example \(\max(R^2,r^2)/(1+R^2)\) is between
\(\min(1,r_-^2)/2\) and \(\max(1,r_+^2)\). It follows that (4.2) is
bounded. Fubini with the locally integrable logarithmic kernel gives
\(\omega+dd^c_Cw=\sigma\) distributionally and the bounded psh
representative, so \(w\in E^1\) and \(H(w)=+\infty\).

The total function of the stationary ray is independent of \(\zeta\), is
bounded, and has homogeneous Monge--Ampère measure zero in dimension
two. Comparison on each finite annulus gives the actual weak geodesic.
It remains bounded on every modification, so every total-space generic
coefficient is zero. Thus the datum is zero. The trivial model measure
is \(\delta_{v_{\rm triv}}\) and its entropy is zero. The remaining
statements, including the product extension, follow from Section 3's
product calculation. \(\square\)

This rules out a formula for all \(E^1\) starting potentials with no
entropy hypothesis. It does not refute a theorem for finite-entropy
\(E^1\) starts. Section 8 gives a sufficient transfer mechanism covering
such starts, with its unproved global inputs expressly retained.

\section{5. Exponential floors on an arbitrary compact Kähler
manifold}\label{exponential-floors-on-an-arbitrary-compact-kuxe4hler-manifold}

\textbf{Theorem 5.1.} Let \(u_t\) be an actual bounded weak ray from
zero on \((X,\omega)\), with \(u_t\leq0\). Suppose a fixed
\(\omega\)-psh function \(\psi\) satisfies \[
 u_t\geq\psi\quad(t\geq0),\qquad
 C_b:=\int_Xe^{-b\psi}\,d\mu_0<\infty\quad\text{for every }b>0.
                                                               \tag{5.1}
\] Write \(E(u_t)=-ct\). Then \(c\geq0\). If \(c=0\), the ray is zero.
If \(c>0\), \(H(u_t)/t\to+\infty\) and its actual ordinary datum is
zero.

\textbf{Proof.} By (1.7) with \(v=0\),
\(\int(-u_t)\operatorname{MA}(u_t)\geq ct\). For any probability \(\nu\)
and reference probability \(\mu\), and every bounded measurable \(f\),
\[
 \operatorname{Ent}(\nu\mid\mu)\geq\int f\,d\nu-\log\int e^f\,d\mu.
                                                               \tag{5.2}
\] If entropy is finite, Jensen with the probability \(\nu\) gives this
inequality by applying concavity of log to \(e^f/(d\nu/d\mu)\) on its
positive density set. If entropy is infinite it is automatic; singular
measures also have infinite entropy by definition. Applying (5.2) to the
bounded function \(f=-bu_t\) gives \[
 H(u_t)\geq bct-\log C_b.                                  \tag{5.3}
\] First take \(t\to\infty\) at each fixed \(b\), and then take
\(b\to\infty\). This proves the infinite limit for \(c>0\) without
uniformity in \(b\) and without claiming finite entropy at every time.

The total function \(U(x,\zeta)=u_{-\log|\zeta|^2}(x)\) is locally
bounded above and psh after adding a smooth local potential. It extends
across the central divisor by removability. The lower bound by \(\psi\)
holds off that divisor, which suffices for Fubini: \[
 \int_{X\times\Delta}e^{-bU}\,d\mu_0\,dA(\zeta)<\infty
                 \quad\text{for every }b>0.                 \tag{5.4}
\] On a smooth modification, let a prime divisor be \(\{w_1=0\}\) at a
generic point, and let the Jacobian order be \(k\). A positive
coefficient \(\lambda\) would give \(U\circ\pi\leq\lambda\log|w_1|^2+C\)
by divisorial current subtraction. Change of variables in (5.4) would
then require integrability of \(|w_1|^{-2b\lambda+2k}\), which fails for
\(b\lambda\geq k+1\). Since every \(b>0\) is allowed, all generic
divisorial coefficients are zero, including those over the time divisor.
The ordinary datum consequently vanishes on every divisorial valuation;
divisorial determination in the NA psh space identifies it with zero (or
this last statement may be read directly as vanishing of its ordinary
divisorial datum).

If \(c=0\), every term in the sum (1.2) is nonpositive. Its sum zero
forces the \(j=0\) term \(\int_Xu_t\omega^n\) to vanish. Thus \(u_t=0\)
almost everywhere. Equality of psh representatives from
almost-everywhere equality implies \(u_t=0\) everywhere. \(\square\)

The theorem is a sufficient analytic criterion. It does not assert that
all nonmaximal rays have floors of this kind, or independently
reconstruct the general maximalization theorem. The previously supplied
polar counterexample is consistent with it but is not needed in its
proof.

\section{6. The SNC log-canonical defect, including nonreduced
fibers}\label{the-snc-log-canonical-defect-including-nonreduced-fibers}

\textbf{Proposition 6.1.} Let \(p:\mathcal Y\to\mathcal X\) be a
modification between smooth manifolds, an isomorphism off a divisor
\(B\) on \(\mathcal X\), with \(B_{\rm red}\) SNC. Put
\(B_{\mathcal Y}=p^*B\) and suppose its reduced support is SNC. Define
\[
 C_{\mathcal X}=K_{\mathcal X/(X\times\mathbf P^1)}+B_{\rm red}-B,
\] and \(C_{\mathcal Y}\) similarly. Then \[
 C_{\mathcal Y}=p^*C_{\mathcal X}+\Delta,
 \quad\Delta=K_{\mathcal Y/\mathcal X}
       +(B_{\mathcal Y})_{\rm red}-p^*B_{\rm red},
 \quad 0\leq\Delta\leq K_{\mathcal Y/\mathcal X}.           \tag{6.1}
\]

\textbf{Proof.} The identity follows from \(B_{\mathcal Y}=p^*B\) and
transitivity of canonical divisors. Strict transforms have coefficient
zero in \(\Delta\). At an exceptional prime \(F\), its center lies in
\(B\); set \(a=\operatorname{ord}_FK_{\mathcal Y/\mathcal X}\) and
\(m=\operatorname{ord}_Fp^*B_{\rm red}\geq1\). The coefficient is
\(a+1-m\). We give the local proof of its nonnegativity. Near a smooth
point of the center write \(B_{\rm red}=\{z_1\cdots z_r=0\}\) and
consider
\(dz_1/z_1\wedge\cdots\wedge dz_r/z_r\wedge dz_{r+1}\wedge\cdots\wedge dz_N\).
At the generic point of \(F\), the pullback of each logarithmic
differential has at most a \(dw_1/w_1\) pole plus a holomorphic term. In
their exterior product at most one such pole can occur, since
\(dw_1\wedge dw_1=0\). Its order is consequently at least \(-1\). But
that order is \(a-m\), so \(a+1-m\geq0\). The bound above by \(a\)
follows from \(m\geq1\). There are no horizontal exceptional divisors
under the hypothesis. This proves (6.1). \(\square\)

If \(\beta_k\) are Kähler classes on such \(\mathcal Y_k\) and
\(\beta_k^n\cdot K_{\mathcal Y_k/\mathcal X}\to0\), positivity gives \[
 0\leq\beta_k^n\cdot\Delta_k
     \leq\beta_k^n\cdot K_{\mathcal Y_k/\mathcal X}\to0.     \tag{6.2}
\] This implication is valid. It does not follow from convergence of
\(\beta_k^{n+1}\) alone. The multiplicities of \(B\) have not been
replaced by one in either \(C\); omitting \(B_{\rm red}-B\) destroys
(1.5).

\section{7. Exactly what model recovery would
prove}\label{exactly-what-model-recovery-would-prove}

We state a conditional result in the genuine compact-Kähler NA theory.
Its extra hypothesis is an explicit input, not a solution of the
original atom.

The external positive-product formulas used in this section are {[}MPWN,
arXiv:2509.09442v1, Propositions 9.1.3 and 9.1.5{]}: on a smooth
dominating SNC model and for a real vertical \(D\) with \(f_D\geq1\),
writing \(\Gamma=A+D\), \(\phi=P_\alpha(f_D)\), they assert \[
 E^{\rm NA}(\phi)=\frac{\operatorname{Vol}(\Gamma)}{(n+1)V},\qquad
 H^{\rm NA}(\phi)=V^{-1}\langle\Gamma^n\rangle\cdot C_{\mathcal X}.
                                                               \tag{7.1}
\] Here the underlying mass formula, their Proposition 4.2.2, is
\(\operatorname{MA}^{\rm NA}(\phi)=V^{-1}\sum b_i\langle\Gamma^n\rangle\cdot E_i\,\delta_{v_{E_i}}\).
The coefficient calculation (1.5) proves the entropy formula from that
mass formula. Their energy formula follows from their mixed
positive-product formula. The source constructs relatively Kähler
approximants increasing uniformly to the continuous envelope, and
identifies their intersections with positive products. Thus (7.1) is
used for PL obstacles in its precise fixed-smooth-\(X\) scope; it is not
entropy convergence for arbitrary rays. Central translations allow
\(f_D\geq1\) and are undone at the end.

\textbf{Theorem 7.1 (model recovery from specified special data).} Fix
such \(\mathcal X,D,\Gamma,\phi\). Suppose actual smooth vertical SNC
modifications and effective vertical defects have been supplied: \[
 p_k:\mathcal Y_k\to\mathcal X,\quad E_k\geq0,\quad
 \beta_k=p_k^*\Gamma-E_k\text{ Kähler}.
\] Assume the modifications and defects are \(\mathbf C^*\)-equivariant,
and assume the three exact convergence statements \[
 \begin{split}
 \beta_k^{n+1}&\longrightarrow\operatorname{Vol}(\Gamma),\\
 \beta_k^n\cdot K_{\mathcal Y_k/\mathcal X}&\longrightarrow0,\\
 \beta_k^n\cdot p_k^*\xi&\longrightarrow\langle\Gamma^n\rangle\cdot\xi
       \quad\text{for every fixed smooth real closed $(1,1)$ class }\xi.
 \end{split}                                                  \tag{7.2}
\] Then the actual relatively Kähler model metrics
\(\phi_k=f_{p_k^*D-E_k}\) satisfy \(\phi_k\leq\phi\),
\(E^{\rm NA}(\phi_k)\to E^{\rm NA}(\phi)\), and
\(H^{\rm NA}(\phi_k)\to H^{\rm NA}(\phi)\). They converge strongly,
using the strong quasi-metric and fixed-output integration facts stated
below.

\textbf{Proof.} Effectivity gives \(\phi_k\leq f_D\) and \(\phi_k\) is
psh, so it lies below \(P_\alpha(f_D)\). Model energy and (7.1) give
energy convergence from the first line of (7.2). Model entropy,
Proposition 6.1 and the other two lines give \[
 H^{\rm NA}(\phi_k)=V^{-1}\beta_k^n\cdot(p_k^*C_{\mathcal X}+\Delta_k)
     \longrightarrow V^{-1}\langle\Gamma^n\rangle\cdot C_{\mathcal X}
     =H^{\rm NA}(\phi).                                     \tag{7.3}
\] Every fixed real class has a smooth representative, and can be
written as a difference of positive classes by adding a sufficiently
large Kähler class. Thus the third line applies to \(C_{\mathcal X}\)
even when its representative is not positive.

For strong convergence, the source strong theory uses the
translation-invariant gap \[
 J(a,b)=E^{\rm NA}(a)-E^{\rm NA}(b)+\int(b-a)\operatorname{MA}^{\rm NA}(a).
\] For \(b=\phi_k\leq a=\phi\), its nonnegativity and the energy
convergence give
\(0\leq J(\phi,\phi_k)\leq E^{\rm NA}(\phi)-E^{\rm NA}(\phi_k)\to0\).
After normalizing suprema to zero, the strong quasi-metric result
{[}MPWN, Theorem 5.2.2{]} gives strong convergence. Write
\(c_k=\sup\phi_k\), \(c=\sup\phi\). The normalized energy convergence
and the unnormalized energy convergence then imply \(c_k\to c\), since
\(E^{\rm NA}(\phi_k-c_k)=E^{\rm NA}(\phi_k)-c_k\). Restoring constants
proves unnormalized strong convergence. \(\square\)

The abstract three convergence assumptions in (7.2) are not counted as
established globally here. Trusiani's September preprint {[}T26, Theorem
B and Lemmas 3.4--3.5{]} asserts special Kähler Fujita approximations,
and its Section 5 asserts that for a big test configuration their
defects and centers can be made vertical and \(\mathbf C^*\)-invariant.
This is the intended route to (7.2) in dimension \(n+1\). The source's
proof of its discrepancy input has the flaw in Section 10; in this paper
that chain is not promoted to an independently verified unconditional
theorem.

MPWN's constructive proof of Theorem 9.2.1 gives an additional two-stage
recovery: retract an MA measure to SNC dual complexes with strong and
discrepancy-integral convergence (Corollary 7.2.3); solve their NA MA
equations by Theorem B and their continuity Theorem 8.2.1; approximate
the retraction-invariant continuous obstacle by PL functions on
subdivisions of the \emph{same} dual complex. The discrepancy is
continuous and PL on that fixed complex, so weak measure convergence
there gives entropy convergence. An arbitrary new model with new
exceptional valuations would not have that justification. Combining this
procedure with a validated Theorem 7.1 input would give a diagonal
relatively Kähler model recovery on \(E^{1,\rm NA}\) with finite
entropy. That final unconditional model recovery is not asserted here.

\section{8. Fixed-output integration and transfer to actual
rays}\label{fixed-output-integration-and-transfer-to-actual-rays}

This section isolates a logically sufficient mechanism, including
finite-entropy \(E^1\) starting points. It explicitly avoids treating
the maximum of two rays as a geodesic.

The NA background facts used here are {[}MPWN, Lemma 5.0.5 and Theorem
5.2.2{]}: strong convergence means weak convergence on divisorial
valuations together with energy convergence; \(J(\phi_j)\) is bounded
for a strongly convergent sequence; and if \(\phi_j\) converges weakly
with bounded \(J\), then \[
 \int|\phi_j-\phi|\,\operatorname{MA}^{\rm NA}(\psi)\to0
       \quad(\psi\in E^{1,\rm NA}).                         \tag{8.1}
\] Here is the check on the relevant proof. Normalize the suprema,
approximate \(\psi\) decreasingly by model \(\psi_k\), and write
\(\mu_k=\operatorname{MA}^{\rm NA}(\psi_k)\). At fixed \(k\), their
finite model support makes integrals converge by divisorial convergence.
The NA energy-pairing estimate (MPWN Proposition 5.0.1/Corollary 5.0.2)
bounds both errors with \(\mu-\mu_k\) uniformly in \(j\), by a constant
times \(J(\psi_k,\psi)^{2^{-n}}\) and a uniformly bounded energy factor.
This tends to zero. Thus signed integrals converge. Apply the same
argument to \(\widetilde\phi_j=\max(\phi_j,\phi)\), which converges
weakly and has bounded \(J\) by normalized order and energy
monotonicity. Finally use \[
 |\phi_j-\phi|=2(\widetilde\phi_j-\phi)-(\phi_j-\phi).
                                                               \tag{8.2}
\] This is the required minus sign; the plus sign in the supplied text's
final identity is a typographical error. The proof establishes (8.1),
not arbitrary entropy continuity.

\textbf{Proposition 8.1 (common upper actual ray).} Let
\(\eta_j\to\phi\) strongly in the preceding NA theory. Assume that the
associated \emph{actual} maximal rays from a fixed \(w\in E^1\) exist,
preserve order in the NA datum, and obey \[
 E(r_t^{w,\theta})=E(w)+tE^{\rm NA}(\theta).                \tag{8.3}
\] Then \(d_1(r_t^{w,\eta_j},r_t^{w,\phi})\to0\), uniformly on bounded
time intervals. The assertion is conditional on the calibrated
actual-ray correspondence (8.3), not on an unspecified roof isometry.

\textbf{Proof.} Let \(\lambda_j=\max(\eta_j,\phi)\) and let \(s_j\) be
its maximal ray from \(w\). The NA supporting inequality and (8.1) give
\[
 0\leq E^{\rm NA}(\lambda_j)-E^{\rm NA}(\phi)
 \leq\int(\lambda_j-\phi)\operatorname{MA}^{\rm NA}(\phi)
 \leq\int|\eta_j-\phi|\operatorname{MA}^{\rm NA}(\phi)\to0.
\] Order preservation implies
\(s_{j,t}\geq r_t^{w,\eta_j},r_t^{w,\phi}\). The triangle and ordered
metric identities, followed by (8.3), yield \[
 d_1(r_t^{w,\eta_j},r_t^{w,\phi})
 \leq t\{2E^{\rm NA}(\lambda_j)-E^{\rm NA}(\eta_j)-E^{\rm NA}(\phi)\}\to0.
                                                               \tag{8.4}
\] The same estimate is uniform for \(0\leq t\leq T\). \(\square\)

The precise source of the proposed correspondence is {[}MP, Theorem
5.1.7{]}: for each \(w\in E^1\) and each \(\phi\in E^{1,\rm NA}\) it
constructs a unique maximal geodesic ray from \(w\) with datum \(\phi\).
Its definition of maximality implies the order comparison. Its proof
takes decreasing smooth starts and decreasing model data, then uses
energy affinity and decreasing energy continuity. For squared-log
convention (1.1), its model slope identities must first be calibrated as
in Section 2. Proposition 8.1 states that requirement explicitly rather
than granting a mismatched formula.

Here is a transfer result independent of an entropy continuity claim.
Let \(G_C=M_C-H\), and write \(G^{\rm NA}\) for its calibrated
mixed-energy slope functional.

\textbf{Theorem 8.2 (a sufficient entropy upper transfer).} Fix \(\phi\)
with finite \(H^{\rm NA}(\phi)\) and a finite-entropy \(w\in E^1\).
Suppose the actual maximal rays exist, preserve order in both their
starting potentials and their NA data, and satisfy (8.3). Suppose model
data \(\eta_j\) satisfy strong and entropy convergence to \(\phi\), with
\(G^{\rm NA}(\eta_j)\to G^{\rm NA}(\phi)\). Assume also the following
two calibrated analytic inputs:

\begin{enumerate}
\def\labelenumi{\arabic{enumi}.}
\tightlist
\item
  From each smooth starting potential, the actual ray for \(\eta_j\) has
  \(M_C\) slope
  \(M^{\rm NA}(\eta_j)=H^{\rm NA}(\eta_j)+G^{\rm NA}(\eta_j)\).
\item
  The target maximal ray has finite \(G_C\) slope equal to
  \(G^{\rm NA}(\phi)\).
\end{enumerate}

Then the maximal ray from \(w\) satisfies \[
 \limsup_{t\to\infty}\frac{H(r_t^{w,\phi})}{t}
      \leq H^{\rm NA}(\phi).                               \tag{8.5}
\] This is a proved implication with explicit additional hypotheses. It
does not establish those inputs on all of the original residual domain.

\textbf{Proof.} Choose smooth \(w_j\) with \(d_1(w_j,w)\to0\) and
\(H(w_j)\to H(w)\) using {[}BDL, Theorem 3.2{]}. Continuity of the
energy terms then gives \(M_C(w_j)\to M_C(w)<\infty\). We verify that
changing the start causes no hidden time-dependent error. If two actual
maximal rays have the same datum but starts \(a,b\), compare both to the
maximal ray from \(\max(a,b)\); order in the start follows from the
defining maximality. By (8.3), the ordered metric bound is independent
of time: \[
 d_1(r_t^{a,\theta},r_t^{b,\theta})
    \leq 2E(\max(a,b))-E(a)-E(b).                           \tag{8.6}
\] The right-hand side tends to zero when \(a=w_j\to b=w\) strongly. To
check this fact directly, let \(p_j=P(w_j,w)\). The rooftop metric
identity and monotonicity show \(E(w)-E(p_j)\to0\) and
\(E(w_j)-E(p_j)\to0\). More directly, {[}D, arXiv:1409.2072v2, Theorem
5.5{]} gives \(\int|w_j-w|\operatorname{MA}(w)\leq C d_1(w_j,w)\to0\).
The supporting inequality (1.7) therefore gives
\(0\leq E(\max(w_j,w))-E(w)\leq\int(w_j-w)_+\operatorname{MA}(w)\to0\).
This proves the required maximum continuity; it is not a statement about
entropy.

Now set \(r_j=r^{w_j,\eta_j}\). Combine (8.6) with (8.4) for start \(w\)
to get \(d_1(r_{j,t},r_t^{w,\phi})\to0\) for every fixed \(t\).
Convexity and the finite model slope give \[
 M_C(r_{j,t})\leq M_C(w_j)+tM^{\rm NA}(\eta_j).
\] Fix \(t\) before taking \(j\to\infty\). Lower semicontinuity {[}BDL,
Theorem 4.7{]} gives \[
 M_C(r_t^{w,\phi})
   \leq M_C(w)+t\{H^{\rm NA}(\phi)+G^{\rm NA}(\phi)\}.
                                                               \tag{8.7}
\] Only after this inequality has been established for every fixed time
do we divide by \(t\) and take \(t\to\infty\). Subtract the assumed
finite matching \(G_C\) slope; the fixed finite constant \(M_C(w)\)
disappears. This proves (8.5). \(\square\)

For \(w=0\) this proof reduces exactly to the fixed-time
lower-semicontinuity argument proposed in Own178, once its global inputs
and normalization are corrected. Finite-entropy \(E^1\) starts can
therefore be accommodated by this same transfer mechanism.
Infinite-entropy starts cannot be inserted into (8.7): its starting
constant would be infinite, as Theorem 4.1 makes concrete.

\section{9. The model log-volume exponent and the entropy lower
bound}\label{the-model-log-volume-exponent-and-the-entropy-lower-bound}

This section verifies the exponent that an ordinary entropy lower-bound
argument must use. It does not assume that the lower bound's stated
proof in {[}MP{]} already has compatible factors.

\textbf{Proposition 9.1 (SNC integrability with multiplicities).} Let
\(\rho:\mathcal X\to X\times\Delta\) be a smooth dominating modification
with SNC central support, \(B=\sum b_iE_i\), and relative canonical
divisor \(K_{\mathcal X/(X\times\Delta)}=\sum a_iE_i\). Put
\(c_i=a_i+1-b_i\). Let \(F_t\) be a smooth-on-the-generic-fiber model
log-volume weight such that locally \[
 F\circ\rho=-\sum c_i\log|z_i|^2+O(1),                    \tag{9.1}
\] where the bounded remainder is smooth on the model. Then for large
\(t=-\log|\zeta|^2\), \[
 \int_Xe^{F_t}\,d\mu_0\leq C(1+t)^n.                     \tag{9.2}
\]

\textbf{Proof.} Near a point where \(r\) central components meet, choose
coordinates with \(\zeta=\prod_{i=1}^rz_i^{b_i}\). A holomorphic unit in
this expression is absorbed in one coordinate by taking a local root.
The pullback of a smooth positive total volume from \(X\times\Delta\) is
comparable to \(\prod|z_i|^{2a_i}\) times Euclidean volume, since the
Jacobian divisor is the stated relative canonical divisor.
Multiplication by \(e^F\) therefore gives the radial factor \[
 \prod|z_i|^{2(a_i-c_i)}=\prod|z_i|^{2(b_i-1)}.
                                                               \tag{9.3}
\] This is the exact cancellation; replacing \(c_i\) by \(a_i\) when
\(b_i>1\) would give a different measure.

Write \(q_i=-\log|z_i|^2\). Euclidean radial area contributes
\(e^{-q_i}dq_i\), so (9.3) contributes \(\prod e^{-b_iq_i}dq_i\). The
base satisfies \(t=\sum b_iq_i\), and its Euclidean area contributes
\(e^{-t}dt\). The fiber density after division by base area is bounded
by a constant times the volume of the simplex slice \(\sum b_iq_i=t\),
\(q_i\geq L_i\), namely \(O((1+t)^{r-1})\). Angular integration has a
finite multiplicity depending on the fixed integers \(b_i\), and the
remaining coordinates range in a bounded polydisk. This can also be
obtained by convolution of the \(r\) exponential radial densities; the
resulting base density is \(e^{-t}\) times a polynomial bound of degree
\(r-1\). A finite coordinate cover and partition of unity gives (9.2),
because \(r-1\leq n\). Smooth weights and coordinate units change only
the comparison constants. \(\square\)

The smooth product pair with its central fiber is log canonical: the
logarithmic-form argument of Proposition 6.1 gives \(c_i\geq0\). Thus
(9.1) also bounds \(F_t\) below uniformly. Consequently
\(\log\int e^{F_t}\mu_0=O(\log(1+t))\), in particular the upper estimate
required for entropy duality. For a maximal ray for which the calibrated
fixed-output mixed slope gives \[
 \lim_{t\to\infty}\frac1t\int F_t\operatorname{MA}(r_t)
       =\int f_C\,\operatorname{MA}^{\rm NA}(\phi),          \tag{9.4}
\] (5.2) and (9.2) prove the corresponding model entropy lower bound.
Taking the supremum over SNC retractions gives the full NA entropy lower
bound if (9.4) is available on that ray. This last mixed-slope input is
exactly where {[}MP, Theorem 5.3.4{]} enters its Proposition 6.2.2.
Equation (9.4) is not established anew on all rays in this paper.

In the logarithmic convention of (2.5), the compatible weight is half of
\(F\) and one must use \(e^{2F_L}\) in entropy duality. The paired
output must consequently be multiplied by two when probability entropy
is used. BHJ's \(e^{2\Psi}\) convention reflects this. Own178's use of
squared-log divisor coefficients, unsquared time, and unchanged
probability entropy cannot simply import an \(e^{F_L}\) bound and a unit
slope simultaneously.

\section{10. A local defect in the imported analytic discrepancy
proof}\label{a-local-defect-in-the-imported-analytic-discrepancy-proof}

Theorem B of {[}T26{]} invokes its Proposition 4.1, which in turn
imports the analytic proof of {[}Tproj, Proposition 3.1{]}. We inspected
that proof, including its determinant estimate. There is a concrete
false geometric assertion in the step making its residual form Kähler
near a generic exceptional point.

\textbf{Proposition 10.1.} A normalized blow-up of an ideal on a smooth
compact Kähler surface need not have a differential injective in the
normal direction of an exceptional divisor. In the example below, the
residual form used by the cited analytic proof has zero normal value
there, for every multiple of the background form.

\textbf{Proof.} In a bidisk around a point on
\(\mathbf P^1\times\mathbf P^1\), take \(I=(x^2,y^3)\). This defines a
global coherent ideal supported at that point by using \(\mathcal O\)
off the point; on the punctured overlap it is the unit ideal. Near the
generic exceptional divisor of its normalized blow-up, a smooth toric
chart has coordinates \((w,v)\) with \(v\ne0\) and \[
 x=w^3v,\qquad y=w^2v.                                    \tag{10.1}
\] Here the primitive valuation is
\((\operatorname{ord}_Ex,\operatorname{ord}_Ey)=(3,2)\). To verify the
chart without appealing to a diagram, start with the blow-up chart where
the ratio \(T=y^3/x^2\) is nonzero. Its function field has \(v=T\) and
\(w=x/y\). The relations give \(y=w^2v\), \(x=w^3v\); \(w\) is integral
over the chart, since \(w^2=y/v\). The resulting ring
\(\mathbf C[w,v,v^{-1}]\) is normal, finite and birational over that
chart. It is its normalization. Thus \(w=0\) is a genuine generic
exceptional divisor on a normalized blow-up.

Along \(E=\{w=0\}\) both \(\partial x/\partial w=3w^2v\) and
\(\partial y/\partial w=2wv\) vanish. In fact the entire differential is
zero there. No pullback of a smooth background Kähler form can be
positive on \(\partial_w\) at \(E\).

The ideal norm gives exactly \[
 \log(|x|^4+|y|^6)
 =6\log|w|^2+\log(|v|^4+|v|^6).                           \tag{10.2}
\] After subtracting \(6[E]\), the local curvature of this potential is
\(dd^c_C\log(|v|^4+|v|^6)\), which has no normal component. Adding any
\(C\rho^*\omega\) does not change that fact. One can choose the global
partition of unity in the cited construction to have one chart weight
identically one near the point, making (10.2) precisely its local
residual potential. More generally smooth generator changes or partition
weights contribute terms through \((x,y)\); their first and mixed normal
second derivatives at \(w=0\) vanish as well. Thus the residual form
cannot be made Kähler there by increasing the background coefficient.

For comparison the ordinary Jacobian determinant in (10.1) is \(w^4v\),
so \(a_E=4\). The identity (10.2) has ideal order \(6\). The cited proof
needs nonvanishing of the full residual determinant to cancel it in its
inequality (3.4); here that determinant vanishes. \(\square\)

This is a counterexample to a proof step, \textbf{not} to the
discrepancy estimate
\(a_E\leq N(\operatorname{ord}_EI-\operatorname{ord}_E\mathcal J(I))\)
or to Theorem B. There is also an algebraic proof in {[}Tproj, Appendix
A, Proposition A.1{]}. That argument uses relative Nadel vanishing to
generate \(\mathcal O(K-ND)\) at a generic Rees divisor and the
subadditivity inclusion \(\mathcal J(I^N)\subset\mathcal J(I)^N\). Its
printed domain is smooth algebraic varieties. A repair for analytic
ideals on arbitrary compact Kähler manifolds would require an exact
analytic relative vanishing statement and its application to this
generation argument, or a different local analytic proof. We have not
supplied such a repair. Importing the appendix with its algebraic
hypotheses suppressed would not satisfy the assignment's
source-application rules.

\section{11. Audit of Own178 and the precise unresolved original
scope}\label{audit-of-own178-and-the-precise-unresolved-original-scope}

The following audit concerns actual assertions, not its author's
confidence or an earlier review result.

\begin{longtable}[]{@{}
  >{\raggedright\arraybackslash}p{(\columnwidth - 2\tabcolsep) * \real{0.5000}}
  >{\raggedright\arraybackslash}p{(\columnwidth - 2\tabcolsep) * \real{0.5000}}@{}}
\toprule\noalign{}
\begin{minipage}[b]{\linewidth}\raggedright
Step in Own178
\end{minipage} & \begin{minipage}[b]{\linewidth}\raggedright
Outcome of this audit
\end{minipage} \\
\midrule\noalign{}
\endhead
\bottomrule\noalign{}
\endlastfoot
\(dd^c=i/(2\pi)\), time \(-\log|\zeta|\), squared-log coefficients, unit
energy slope & Refuted as a coherent combination by Proposition 2.1. Use
(1.1), or (2.5)--(2.6) with half entropy. \\
\(0\leq\Delta\leq K_{\rm relative}\) & Proved in Proposition 6.1 with
the nonreduced correction retained. \\
Fixed-form convergence plus canonical-defect control implies model
entropy recovery & Proved in Theorem 7.1 under the exact supplied-data
hypotheses (7.2); volume alone is insufficient. \\
Ordered model energy convergence implies strong convergence & Valid
using the explicitly stated NA strong theory; removal and restoration of
suprema is necessary. \\
Two-stage MPWN obstacle recovery & Source-supported within its fixed
smooth Kähler/dual-complex scope; it does not by itself give relatively
Kähler model entropy recovery. \\
Fixed-output integration and common upper ray & The scalar sign in the
supplied Lemma 5.0.5 text needs correction (8.2). The transfer estimate
(8.4) is valid under the calibrated actual-ray identities. \\
Fixed-time Mabuchi lower semicontinuity before the time limit & Valid
sufficient mechanism, Theorem 8.2. No entropy continuity along arbitrary
strong sequences is used. \\
Model Mabuchi/mixed slope imported from SD and MP & Must be normalized
anew; SD's printed Green--Riesz factor fails the constant-model test.
This paper does not certify the whole source chain for all rays. \\
Finite-Mabuchi-slope maximality & MP Proposition 6.3.1 states the
compact-Kähler transfer; its proof refers to Li's local integrability
argument. The local structure is plausible and consistent with
squared-log time, but a full transported proof on every starting-point
scope is not submitted here. \\
All data of an arbitrary \(E^1\) ray have finite NA energy & Proposed
maximalization/truncation proof requires the actual dominance and energy
identities with their precise normalization; it is not established as an
unconditional new theorem here. \\
Global maximal entropy equality and arbitrary-ray dichotomy & Not
established by this audit. The preceding partial results do not close
P19.1. \\
\end{longtable}

For finite-entropy starts, the proof of Theorem 8.2 shows concretely how
starting-point approximation can preserve the initial finite Mabuchi
constant. It is therefore not correct to list ``smooth start versus
finite-entropy start'' as an inevitable obstruction to this mechanism.
However its global calibrated model recovery and mixed-slope hypotheses
still need to be established. For infinite-entropy starts, Theorem 4.1
gives an actual bounded stationary obstruction, so no formula may
discard that case without an additional term or a restriction.

Li's proof of nonmaximality implying infinite entropy can be isolated as
follows. If a ray and its maximalization from the same start have
positive affine energy gap \(ct\), and if
\(\int_{X\times\Delta}\exp(b(\widehat U-U))<\infty\) for every \(b>0\),
then Fubini in squared-log time yields a sequence \(t_j\to\infty\) on
which \(\int_Xe^{b(\widehat u_{t_j}-u_{t_j})}\mu_0\leq e^{t_j}\).
Entropy duality, truncation of the nonnegative test function, and (1.7)
give \(H(u_{t_j})\geq bct_j-t_j\). If the entropy slope exists, allowing
\(b\) to increase forces it to be infinite. This is the analytic
deduction in {[}Li, Lemma 3.1 and proof of Theorem 1.2{]}; the
all-exponential integrability assertion for arbitrary transcendental
rays remains an imported prerequisite if this deduction is used
globally. The floor theorem of Section 5 has its own directly proved
integrability and does not use that import.

The residual catalogue calls for the entropy-slope statement for
arbitrary finite-energy weak rays outside the already known smooth
polarized maximal-ray case. This paper establishes neither all maximal
finite-NA-entropy recovery nor the complete arbitrary-ray classification
on that domain. The model theorem in Section 3 is a special case, and
the conditional transfer theorems cannot be counted as a full atom by
the completion criterion. No fixed-\(X\) result here supplies the finer
singularity preservation of P19.2, the genuinely degenerating-family
compactness of P19.3, or the equivariant/singular metric signatures of
P19.4.

Concrete remaining requirements are: (i) repair or replace the analytic
Rees-discrepancy proof underlying the September special-Fujita input,
with all analytic hypotheses checked; (ii) finish the normalized model
ray and mixed-slope source audit and then establish the unconditional
relatively Kähler energy/entropy recovery in the full NA domain; (iii)
transport the general maximalization and nonmaximal integrability proof,
including the entropy-slope existence argument and finite-entropy
\(E^1\) starting points; (iv) state separately the treatment of
infinite-entropy starts. They are mathematical gaps, not failures of TeX
compilation or of numerical tests.

\section{12. References and application
boundaries}\label{references-and-application-boundaries}

{[}BT{]} E. Bedford and B. A. Taylor, \emph{The Dirichlet problem for a
complex Monge--Ampère equation}, Invent. Math. 37 (1976), 1--44; \emph{A
new capacity for plurisubharmonic functions}, Acta Math. 149 (1982),
1--40. Only bounded local product, comparison and Dirichlet uniqueness
facts are imported here.

{[}Dem{]} J.-P. Demailly, \emph{Complex Analytic and Differential
Geometry}, version of 21 June 2012,
\href{https://people.math.harvard.edu/~demarco/Math274/Demailly_ComplexAnalyticDiffGeom.pdf}{public
primary-text mirror}. Chapter I Theorem 5.24 and its polar-set proof,
Chapter III Section 8 Proposition 8.16 and the ensuing decomposition
proof were inspected in full. The generic-divisor subtraction is applied
to closed positive \((1,1)\) currents on smooth local charts. Chapter
III Section 3 also records the bounded local product and local uniform
convergence facts used above. The author's current hosting endpoints
failed to fetch; the dated primary text at this mirror supplied the
actual proof units, rather than a search-result description.

{[}D{]} T. Darvas, \emph{The Mabuchi geometry of finite energy classes},
Adv. Math. 285 (2015), 182--219,
\href{https://arxiv.org/pdf/1409.2072v2}{arXiv:1409.2072v2}, 1 November
2014. Corollary 4.14 gives the \(E^1\) rooftop metric formula and
Theorem 5.5 gives comparison with
\(I_1(u,v)=\int|u-v|(\operatorname{MA}(u)+\operatorname{MA}(v))\).
Ordered energy distance follows from the rooftop formula; Section 8
derives the needed maximum continuity from that comparison. No ray
isometry theorem is used.

{[}BDL{]} R. J. Berman, T. Darvas and C. H. Lu, \emph{Convexity of the
extended K-energy and the large time behaviour of the weak Calabi flow},
\href{https://arxiv.org/pdf/1510.01260v2}{arXiv:1510.01260v2}, 7
February 2016, Theorems 3.2 and 4.7, and Sections 2 and 4.2. The
pertinent entropy approximation and lower-semicontinuity/convexity proof
units were inspected; their Monge--Ampère existence and finite-energy
continuity foundations remain external. The scale to (1.1) is explicit
in Section 2.

{[}BHJ{]} S. Boucksom, T. Hisamoto and M. Jonsson, \emph{Uniform
K-stability, Duistermaat--Heckman measures and singularities of pairs:
the differential-geometric aspects},
\href{https://arxiv.org/pdf/1603.01026v4}{arXiv:1603.01026v4}, Lemma
3.11 and its local proof. Its logarithmic metric convention uses half
probability entropy and \(e^{2\Psi}\); only its local volume calculation
is compared here, and Section 9 gives an independent derivation with
arbitrary multiplicities.

{[}SD{]} Z. Sjöström Dyrefelt, \emph{K-semistability of cscK manifolds
with transcendental cohomology class},
\href{https://arxiv.org/pdf/1601.07659v2}{arXiv:1601.07659v2}, Section
2.1, Lemma 4.6, Theorem 4.9, Theorem 5.6 and Remark 5.7. Its domain is
smooth compact Kähler \(X\) and smooth dominating relatively Kähler
cohomological configurations, with compatible rays. Its printed \(dd^c\)
and time require the calibration in Section 2. Theorems 4.9 and 5.6 are
not silently imported with their printed unit factor.

{[}MP{]} P. Mesquita--Piccione, \emph{A non-Archimedean theory of
complex spaces and the cscK problem},
\href{https://arxiv.org/pdf/2409.06221v2}{arXiv:2409.06221v2}, 18
September 2025, Definitions 5.1.2 and 5.1.6, Theorems 5.1.4, 5.1.7 and
5.3.4, Propositions 6.2.2 and 6.3.1. Divisorial determination is Theorem
4.3.8. Its log-\(|f|\) compatibility and time convention were inspected;
its ray correspondence is used only with explicit calibration
requirements in Section 8.

{[}MPWN{]} P. Mesquita--Piccione and D. Witt Nyström, \emph{A
transcendental non-Archimedean Calabi--Yau theorem with applications to
the cscK problem},
\href{https://arxiv.org/pdf/2509.09442v1}{arXiv:2509.09442v1}, 11
September 2025, Lemma 5.0.5, Theorem 5.2.2, Corollary 7.2.3,
Propositions 9.1.3 and 9.1.5, constructive proof of Theorem 9.2.1. The
exact hypotheses and formulas used are stated in Sections 7 and 8. Their
fixed-manifold NA foundations are imports; a catalogue description is
not their proof.

{[}Tproj{]} A. Trusiani, \emph{A solution to the Yau--Tian--Donaldson
conjecture through special Fujita approximations}, with an appendix by
S. Boucksom,
\href{https://arxiv.org/pdf/2605.30063v2}{arXiv:2605.30063v2}, 17 August
2026, Proposition 3.1 and Appendix A Proposition A.1. The complete
relevant analytic proof and the algebraic alternative were inspected.
Section 10 isolates a false geometric step without asserting that the
theorem is false. The already known projective claims are not renamed
transcendental results.

{[}T26{]} A. Trusiani, \emph{On the transcendental Yau--Tian--Donaldson
conjecture},
\href{https://arxiv.org/pdf/2609.31089v1}{arXiv:2609.31089v1}, 25
September 2026, Theorem B, Lemmas 3.4--3.5, Proposition 4.1 and Section
5. Its asserted compact-Kähler special-Fujita statement is recorded as a
recent source claim; the unresolved supporting analytic proof
application is expressly retained. Theorem A's trivial-automorphism
restriction is not a restriction in Theorem B, and no stability theorem
from Theorem A is used here.

{[}Li{]} C. Li, \emph{Geodesic rays and stability in the cscK problem},
author
\href{https://sites.math.rutgers.edu/~cl1412/papers/maximal.pdf}{primary
PDF}, corresponding to Ann. Sci. Éc. Norm. Supér. 55 (2022), 1529--1574,
Lemmas 2.35 and 3.1 and proof of Theorem 1.2. The inspected argument
explicitly uses squared-log time. It is discussed as a transfer
prerequisite, not as a projective theorem already proved here in the
transcendental domain.

\section{13. Reproduction and evidence}\label{reproduction-and-evidence}

The staged \texttt{input.json} and \texttt{input-original.json} are
byte-identical. \texttt{assignment-original.txt} is the exact decoded
assignment; \texttt{task-original.json} records the decoded original
task. The preserved source units are mathematical data. Public source
browser responses, with exact-version URLs and retrieval coordinates,
are retained in the workspace; no other campaign or private output was
read. The failed direct PDF retrieval did not cause a change to proxy,
credentials or runtime configuration.

\texttt{source-application-ledger.json} records each imported theorem,
its full application signature, inspected proof units, normalization
transport and unresolved limitations. \texttt{verify.py} checks input
preservation, the exact status schema, the closed-form surface
intersection identities and several numerical values of the explicit
entropy formula (3.4). The numerical checks are diagnostic observations,
not a proof of any theorem. All mathematical claims above are justified
by the displayed proofs and identified external inputs.

Rerun \texttt{python\ verify.py}, then \texttt{bash\ build-paper.sh}.
The build converts this Markdown to standalone \texttt{P19.1-paper.tex}
with the existing Pandoc and invokes the existing \texttt{pdflatex}
twice, saving both actual logs and console diagnostics. Build success
establishes document buildability only. The paper remains an unverified
research draft with PARTIAL status and candidate=false.

\end{document}
